\section{Conclusion}
\label{sec:conclusion}

As the adoption of Retrieval-Augmented Generation (RAG) accelerates, managing the balance between semantic utility and data security has emerged as a fundamental design challenge. In this work, we formalized the semantic reconstruction threat model and empirically demonstrated the inherent trade-off in representation-space sharding: under standard clustering frameworks, maximizing retrieval utility inherently increases the risk of semantic disclosure. 

To navigate this trade-off, we introduced the SHArd Disclosure Estimator (SHADE), a proxy metric for quantifying semantic co-location risk over a distributed topological graph. We subsequently proposed CoShard, a scalable community detection algorithm designed to optimize the SHADE-coherence trade-off. 

Through information-theoretic analysis and synthetic adversarial simulations over domain-specific RAG proxies (MedQA and LegalBench), we demonstrated that CoShard provides a highly competitive, privacy-aware alternative to existing methods. By formally penalizing intra-shard semantic density, CoShard successfully traces a tunable Pareto boundary, actively starving potential adversaries of critical context while maintaining the semantic proximity required by dense retrievers. 

CoShard offers a software-layer approach to enterprise privacy, enabling system administrators to enforce tunable risk mitigation without suffering catastrophic utility degradation. Future work will explore explicit end-to-end generative attacks to further validate these synthetic bounds, and the adaptation of the CoShard objective function for dynamic, streaming RAG architectures.
